\documentclass[10pt]{article}
\providecommand{\FIGDIR}{fig}
\newcommand{\DIGESTDATE}{3 October 2026}
\newcommand{\PERIODSTART}{27 September 2026}
\newcommand{\PERIODEND}{3 October 2026}
\input{preamble}

\begin{document}

% ==================================================================== MASTHEAD
\thispagestyle{empty}
\begin{tikzpicture}[remember picture, overlay]
  \fill[Deep] (current page.north west) rectangle
      ($(current page.north east)+(0,-67mm)$);
  \fill[Teal] ($(current page.north west)+(0,-67mm)$) rectangle
      ($(current page.north east)+(0,-68.4mm)$);
  \foreach \x/\y/\r in {22/12/0.9, 41/26/1.5, 63/9/0.7, 88/31/1.1, 112/17/1.9,
                        138/29/0.8, 159/13/1.3, 178/54/1.0, 196/21/1.6,
                        30/44/0.6, 74/57/1.2, 125/48/0.9, 168/6/0.8, 191/59/0.7,
                        52/18/1.0, 100/60/0.6, 148/20/1.1, 66/49/1.4, 205/45/0.8}
    \fill[Sky, opacity=0.30] ($(current page.north west)+(\x mm,-\y mm)$) circle (\r mm);
\end{tikzpicture}

\vspace*{4mm}
{\sffamily\fontsize{7.6}{9}\selectfont\color{Sky}%
 WEEKLY ROLLUP\;\textbullet\;PhD RESEARCH WATCH\par}
\vspace{2.2mm}
{\sffamily\bfseries\fontsize{27}{29}\selectfont\color{white}%
 Particulate Matter\par}
{\sffamily\fontsize{27}{29}\selectfont\color{Sky}%
 Week in Review\par}
\vspace{5.0mm}
{\sffamily\fontsize{8.4}{10}\selectfont\color{white}%
 \mbox{Week ending \DIGESTDATE}\;\;\textcolor{Teal}{\textbar}\;\;\mbox{\PERIODSTART{} to \PERIODEND}\;\;
 \textcolor{Teal}{\textbar}\;\;\mbox{125 records pooled from 7 daily issues}\par}
\vspace{18mm}

% -------------------------------------------------------------- metric strip
\noindent
\begin{minipage}[t]{0.190\linewidth}\metric{125}{RECORDS POOLED\\(125 UNIQUE DOIs)}{Deep}\end{minipage}\hfill
\begin{minipage}[t]{0.190\linewidth}\metric{42\%}{MEASUREMENT SIDE\\(52 RECORDS)}{Teal}\end{minipage}\hfill
\begin{minipage}[t]{0.190\linewidth}\metric{10}{SUBTOPICS\\ACTIVE}{Sky}\end{minipage}\hfill
\begin{minipage}[t]{0.190\linewidth}\metric{30}{POOLED EFFECT\\ESTIMATES (30 WITH A CI)}{Amber}\end{minipage}\hfill
\begin{minipage}[t]{0.190\linewidth}\metric{28}{ABSTRACT-FREE\\(22\% OF WEEK)}{Clay}\end{minipage}

\vspace{6mm}

% -------------------------------------------------------------- week in one box
\section*{\sffamily\bfseries\color{Deep}The week in one paragraph}
\vspace{-1.5mm}{\color{Teal}\rule{\linewidth}{1.1pt}}\vspace{2.5mm}

\begin{keybox}
{\sffamily\bfseries\small\color{Deep}Volume fell by nearly a quarter and the measurement side took it
all back: 125 records, with \textit{Exposure} and \textit{Sensing} together at 42\,\% after a series-low
33\,\%. The abstract gap did not close --- 28 of 125 (22\,\%) --- and this rollup is being written six days
late, because the pipeline stopped for four consecutive runs.}\\[1.4mm]
\footnotesize
\textbf{Scale.} \textbf{125 records} over seven issues on daily totals of \textbf{4}, 15, 22, 25,
\textbf{30}, 15 and 14 --- a 7.5$\times$ range, by far the widest in the series, because Sunday
27 September returned \textbf{four}. Against recent rollups: 125 after 162, 148, 126, 150 and 109. All
\textbf{ten} subtopics are populated for the second week running.

\textbf{The measurement side recovered to 42\,\%.} \textit{Exposure assessment \& modelling} 34
$\rightarrow$ \textbf{32} and \textit{Sensing} 20 $\rightarrow$ \textbf{20} held in absolute terms while
the total fell, so the share reads 54, 40, 44, 33, \textbf{42\,\%} across five weeks. Last week's
series low was a denominator effect, as this watch said at the time it might be; it was not a contraction
of the instrumentation literature.

\textbf{\textit{Occupational \& indoor} 22 $\rightarrow$ 14, and the indoor turn is still a model.}
The band that led the health side last week fell back to joint second with \textit{Burden} (14 each).
The week's indoor records are a school PM$_{10}$ deposit with no abstract, waterpipe caf\'es at
\textbf{285\,$\mu$g\,m$^{-3}$}, a stone-carving dust survey and a traffic-police ocular study --- none of
them an occupied-building intervention with a health endpoint. The AIRVAC protocol that entered the
registry last week is still the only randomised test in view.

\textbf{Abstract-free held at 22\,\%.} \textbf{28 of 125}, against 34 of 162 (21\,\%) --- no improvement,
and the retry pass named as the only fix in four consecutive rollups is \textbf{still not built}.
\textbf{17 of the 28} sit on the measurement side, so the abstract gap continues to fall hardest on
exactly the literature this watch exists to track. The Scholar Gateway connector, which is the
designated recovery route for numbers abstracts omit, has been \textbf{down for the entire period} ---
first with an identity error, now \texttt{ACCESS\_DENIED}.

\textbf{Estimates: 30, every one with an interval.} One estimate per 4.2 records (series 4.4, 2.7, 3.1,
5.2, \textbf{4.2}), and \textbf{30 of 30 carry a CI} --- the second complete week in a row. Two issues
(27 September, 1 October) contributed none.
\end{keybox}

\clearpage

% -------------------------------------------------------------- trend
\section*{\sffamily\bfseries\color{Deep}Subtopic trend and inflection}
\vspace{-1.5mm}{\color{Teal}\rule{\linewidth}{1.1pt}}\vspace{2.5mm}

\noindent\includegraphics[width=\linewidth]{\FIGDIR/w1_subtopic_trend.png}
\figcap{Per-issue subtopic counts across the week from \texttt{state/metrics.csv}, seven issues.
Daily totals 4, 15, 22, 25, 30, 15 and 14 --- a \textbf{7.5$\times$} range. The Sunday minimum
(27 September, four records) is the smallest daily issue in the series; the 30-record Thursday
(1 October) is a full weekday deposit. Five of the week's seven issues were built on a later run than
their entry date, and the last of them --- 3 October --- was built \textbf{six days late}, on
9 October, after four consecutive runs halted during authoring. No gap was left.}

\footnotesize
\textbf{Expanding.} \textit{Burden, policy \& mitigation} \textbf{10 $\rightarrow$ 14} is the week's
largest relative movement, carried by the coking-industry inventory, the Middle East conflict-emissions
analysis, two GBD-type burden records and the Jilin CTM work. \textit{Sensing} holds at \textbf{20} on a
25\,\% smaller week, which is a real gain in share.

\vspace{1.2mm}
\textbf{Contracting.} \textit{Occupational \& indoor} \textbf{22 $\rightarrow$ 14} (above) and
\textit{Mechanistic toxicology} \textbf{20 $\rightarrow$ 13} both give back most of last week's rise;
\textit{Cardiovascular \& metabolic} \textbf{15 $\rightarrow$ 8} halves. \textit{Respiratory} (10
$\rightarrow$ 5) is at a series low. \textit{Reproductive \& developmental} holds at \textbf{8} for the
second week, and \textit{Neuro} (8 $\rightarrow$ 6) and \textit{Other clinical} (8 $\rightarrow$ 5) ease.

\vspace{1.2mm}
\textbf{The inflection worth naming.} \textbf{Last week's black-carbon ordering was not tested, and the
constituent question moved to nitrate instead.} The test this watch set --- whether a non-Chinese cohort
reproduces BC leading the constituent weighting --- went unanswered: no constituent-resolved
cardiovascular cohort appeared this week in any country. What did appear, on the last day of the period,
is the \textbf{China Kadoorie Biobank never-smoker lung-cancer analysis}: per IQR, \textbf{nitrate HR
1.35} (1.23--1.48) ahead of ammonium 1.22 and chloride 1.20, with the mixture at 1.13 and an additive
nitrate\,$\times$\,polygenic-risk interaction. Constituent-specific risk is now reported from three
independent Chinese cohorts in three weeks, each naming a \textit{different} leading constituent ---
BC for stroke and heart failure, nitrate for lung cancer. That is either real outcome specificity or
collinearity in one shared reanalysis product, and the two are not yet distinguishable. \textbf{For a
sensor watch the practical point is unchanged and sharpening}: nitrate is semi-volatile and reads low on
heated and warm-cavity optical sensors, BC is measurable by low-cost absorption, and a mass-only network
resolves neither.

\vspace{2mm}
\noindent\includegraphics[width=\linewidth]{\FIGDIR/f2a_architecture.png}
\figcap{Study architecture over the pooled week. \textbf{\textit{Metadata only} leads again at 28},
ahead of \textit{Modelling / inventory} (21) and \textit{Measurement campaign} (20) --- the second
consecutive week in which the missing-abstract bucket is the largest single group. Primary observational
epidemiology (cohort 15, cross-sectional 8, acute 5, ecological 2) totals \textbf{30}, down from 40.
\textbf{No interventional record entered this week}, against four last week.}

\clearpage
% -------------------------------------------------------------- pooled analytics
\section*{\sffamily\bfseries\color{Deep}Pooled analytics}
\vspace{-1.5mm}{\color{Teal}\rule{\linewidth}{1.1pt}}\vspace{3mm}

\noindent\includegraphics[width=\linewidth]{\FIGDIR/f1_subtopics.png}
\figcap{Pooled subtopic distribution, \textbf{all ten bands populated}. \textit{Exposure} 32 and
\textit{Sensing} 20 put the measurement side at \textbf{52 of 125 (42\,\%)}. On the health side
\textit{Occupational \& indoor} and \textit{Burden} tie at 14, then \textit{Mechanistic toxicology} 13;
\textit{Respiratory \& allergic} at \textbf{5} is the weakest band of the week.}

\vspace{2mm}
\noindent\includegraphics[width=\linewidth]{\FIGDIR/f2b_endpoint.png}
\figcap{Health endpoint over the pooled week, canonicalised. \textbf{``No health endpoint'' pools the
majority} --- 28 abstract-free records plus the exposure, instrument, emissions, optics and
air-quality records. Among human outcomes, respiratory and cardiovascular lead at \textbf{4} each, then
reproductive at 3 and metabolic at 2, with a long tail of single-record labels. A week in which no
outcome label reaches five is a week with no clinical centre of gravity.}

\vspace{2mm}
\noindent\includegraphics[width=\linewidth]{\FIGDIR/f3_metric_geography.png}
\figcap{Left --- particle metric. \textbf{PM$_{2.5}$ alone 39 of 125}, unspeciated PM or emissions 37,
joint PM$_{2.5}$ + PM$_{10}$ \textbf{26} (a high share, driven by the indoor and occupational records
that report both). The tail: \textbf{composition / speciation 7}, size-resolved 4, optical properties 4,
PM$_{10}$ alone 4, bioaerosol 3 and \textbf{ultrafine 1}. Right --- geography. \textit{Global /
multi-region} 31 again overstates reach, pooling reviews, laboratory work and the 18 metadata-only
records whose country is not recoverable. \textbf{China 35} is the largest single-country bar --- 28\,\%
of the week --- then Europe 19, North America 15, East Asia (ex-China) 10 and South Asia 8;
Sub-Saharan Africa fell to \textbf{3}.}

\vspace{2mm}
\noindent\includegraphics[width=\linewidth]{\FIGDIR/f6_lifecourse.png}
\figcap{Life-course windows across the pooled week; counts non-exclusive. \textbf{Preconception /
in utero 11} leads for the first time in the series, carried by the Delhi congenital-heart case-control,
the pregnancy-loss and pregnancy-hypertension cohorts and the Medicaid prenatal autism analysis.
Working age \textbf{8} and older adults \textbf{7} follow; infancy 5 and \textbf{adolescence 2} ---
the weakest window for the sixth consecutive week.}

\clearpage

\noindent\includegraphics[width=\linewidth]{\FIGDIR/f5_forest.png}
\vspace{1mm}
\noindent\includegraphics[width=\linewidth]{\FIGDIR/f5_forest_2.png}
\figcap{All \textbf{30} pooled estimates, ratio and difference scales in separate panels.
\textbf{Not mutually comparable and not meta-analysed}: contrasts include per-IQR, per-SD, per-10 and
per-1\,$\mu$g\,m$^{-3}$, per-year, dichotomised and --- in four cases --- \textbf{unstated} increments,
on HR, OR, RR, IRR and $\beta$ scales; heterogeneity is structural, so \textbf{no $I^2$ is reported} and
none would be meaningful. Readable points: \textbf{the three largest ratios are all from one
case-control} --- the Delhi congenital-heart ORs of \textbf{4.35} (2.41--7.83), 3.36 and 2.82, on
1,115 cases against 441 controls with the PM$_{10}$ increment unstated, which is why they sit so far
right of everything else; \textbf{the tightest intervals are the Medicare smoke-PM mortality pair}
(1.032, 1.027--1.037 and 1.087, 1.076--1.098) on 77.5 million beneficiaries; the four China Kadoorie
constituent estimates share cases and are nested, not independent; and the two-sample Mendelian
randomisation OR for otitis media (1.555, 0.947--2.553) is the only estimate on the panel whose
interval crosses unity by design --- the authors report it as a bounded null. The difference panel
carries the two next-day physical-activity $\beta$ coefficients from the 67-adult UK personal-monitoring
panel ($-$8 minutes per IQR, both PM$_{2.5}$ and PM$_{10}$).}

\vspace{2mm}
\noindent\includegraphics[width=\linewidth]{\FIGDIR/f4_heatmap.png}
\figcap{Subtopic $\times$ architecture, pooled. The densest cells are \textit{Exposure} $\times$
\textit{Metadata only} and \textit{Exposure} $\times$ \textit{Modelling}, then \textit{Sensing} $\times$
\textit{Metadata only} --- the abstract gap is again concentrated on the measurement side, where
\textbf{17 of 52 records (33\,\%)} carry no abstract. The standing finding survives a tenth week: of the
\textbf{52} measurement-side records, \textbf{2 carry any outcome label and none observes a human
outcome} --- one is a modelled carcinogenic risk index, the other an emissions-policy label.}

\vspace{4mm}

% -------------------------------------------------------------- venues + trials
\section*{\sffamily\bfseries\color{Deep}Venues, trial watch, and what the pipeline did}
\vspace{-1.5mm}{\color{Teal}\rule{\linewidth}{1.1pt}}\vspace{2.5mm}

\footnotesize
\textbf{Venues.} A two-way tie at \textbf{9} --- \textit{Environ Sci Technol} and \textit{Atmos
Environ} --- then \textit{Atmosphere} and \textit{Indoor Air} at 7, \textit{Front Public Health} and
\textit{J Environ Sci (China)} at 6, \textit{Atmos Pollut Res} at 5. The leading venue holds 7\,\%.
\textit{Atmos Environ} and \textit{Atmos Pollut Res} remain the main sources of metadata-only records
(Elsevier same-day deposits), exactly as in the last three rollups. \textbf{Preprints supplied 11}, up
from 6 and the highest of the series --- seven of them medRxiv.

\vspace{1.5mm}
\textbf{Recurring authors.} Unchanged: surname frequency is a naming artefact, and \textbf{no
repeat-group cluster is claimed}. No dataset repeated three times this week as CHARLS did last week;
the China Kadoorie Biobank and UK Biobank each supplied one record.

\vspace{1.5mm}
\textbf{Trial watch --- nothing new in the period.} The ClinicalTrials.gov leg was queried on each run
covering the week (\texttt{LastUpdatePostDate} windows 26--29 September, 30 September, 1 October and
2--3 October) and returned \textbf{one hit with nothing kept}: \textbf{NCT07850648} (CAN-EXACT, McGill,
observational COPD-exacerbation cohort, $n$=125, not yet recruiting), whose \texttt{analyze\_endpoints}
shows primaries of exacerbation rate and diagnostic markers with environmental factors only inside a
generic secondary risk-factor measure --- \textbf{no PM or air-quality measurement}, so it was screened
and not added. Carried from the previous period: \textbf{NCT06541691} (COATED-AIR) still reads
\textit{recruiting} past its listed primary completion and still registers \textbf{no exposure or
adherence endpoint}; \textbf{NCT02944682} (HAPIN) posted a registry update. \texttt{trials.json} holds
\textbf{16 trials} over \textbf{27 logged windows}; the 16-vs-24 count discrepancy flagged in the
August monthly was \textbf{audited and closed} on the 2 October run --- it was a drifted tally in
\texttt{counter\_note}, and all counters were re-derived.

\vspace{1.5mm}
\textbf{What the harvesting legs did.} \textbf{Crossref-by-ISSN was the highest-yield leg every
weekday} (40--78 records) and is also the source of most of the week's missing abstracts. \textbf{Europe
PMC} ran clean all week after the 25 September outage (17--75 records/day). The harvester's
\textbf{PubMed sensing query returned 0--3 daily}, and the \textbf{PubMed connector remained
load-bearing}, appending 2--7 records a day that the two harvester queries missed. \textbf{OpenAlex
returned HTTP 429 on every single run of the period} --- it has now contributed nothing for three
weeks and should be treated as dead rather than degraded. \textbf{arXiv} answered (148\,kB) but had no
entry dated inside the window on any day. \textbf{Consensus} returned 10 calibration hits and
\textbf{0 new} records every day, for the sixth consecutive week. \textbf{Scholar Gateway failed on
every attempt}, first \texttt{INVALID\_QUERY: could not resolve user identity}, later
\texttt{ACCESS\_DENIED}; \textbf{it needs the user to reconnect it}. Two real defects were found and
fixed inside the period: the month-boundary gate in \texttt{\_screen\_0929.py} that was discarding
end-of-month preprints on a 1st-of-month issue (5 records recovered), and the batch-free-month regime
in \texttt{plots\_weekly.py} that had shipped a stale legend in the August monthly.

\vspace{2mm}
% -------------------------------------------------------------- gaps
\section*{\sffamily\bfseries\color{Deep}Open gaps}
\vspace{-1.5mm}{\color{Teal}\rule{\linewidth}{1.1pt}}\vspace{2.5mm}

\footnotesize
\begin{itemize}
\item \textbf{The pipeline stopped for four consecutive runs, and this rollup is six days late.} The
6, 7, 8 and 9 October runs harvested and screened but did not author; three of them halted part-way
through writing a single large issue file. The 3 October daily and this weekly were both recovered on
the 9 October run, and authoring is now \textbf{chunked} (\texttt{authoring/<date>/p*.py} plus
\texttt{\_mk\_chunked.py}) so that a partial run leaves usable work behind. This is the largest
availability failure in the series and the defect most worth watching next week.

\item \textbf{The abstract-retry pass, fourth week running.} 22\,\% of the week shipped without an
abstract and \textbf{17 of those 28 are measurement-side records}. The highest-relevance titles are
again among them --- the Central Himalaya CUPI sensor network, a composition-to-volatility
parameterisation, Kuwait PM episode statistics, a size-resolved dry-deposition chamber study. With
\textbf{Scholar Gateway down for the whole period}, the pipeline currently has no working recovery route
for numbers an abstract omits.

\item \textbf{Constituent-specific risk is now reported from three Chinese cohorts naming three
different leading constituents.} Outcome specificity and reanalysis collinearity are not yet
distinguishable, and no non-Chinese constituent cohort has appeared. The test stands.

\item \textbf{None of 52 measurement-side records observes a human outcome}, tenth consecutive week ---
and this week not even one, against one last week.

\item \textbf{Ultrafine and size-resolved number: 5 of 125.} The week's UFP-relevant work is entirely
atmospheric; no health record used a number concentration.

\item \textbf{Pipeline defects carried forward}: no first-author affiliation field; the
\texttt{Cancer}/\texttt{Oncologic} split in \texttt{ENDPOINT\_CANON}; no \texttt{\textbackslash
DIGESTDATE}-vs-\texttt{--date} build guard; the 2026-07-27 corpus missing \texttt{LIFECOURSE};
\texttt{.git} lock debris accumulating on the delete-restricted mount.
\texttt{metrics.csv} continues clean (\texttt{csv.writer}, 541 rows at period close, 0 malformed).
\end{itemize}

\clearpage

% -------------------------------------------------------------- digest
\section*{\sffamily\bfseries\color{Deep}Paper-level digest --- the week by subtopic}
\vspace{-1.5mm}{\color{Teal}\rule{\linewidth}{1.1pt}}\vspace{1mm}

{\fontsize{7.4}{9}\selectfont\color{Slate}Every record that appeared in a daily issue
keeps that issue's critical summary verbatim; nothing here is re-written or re-assessed.
Records surfaced without an abstract are shown as such. Ordered by subtopic, then by the
date of the issue that first carried the record.\par}

\input{digest_body}

% -------------------------------------------------------------- register
\section*{\sffamily\bfseries\color{Deep}Pooled corpus register}
\vspace{-1.5mm}{\color{Teal}\rule{\linewidth}{1.1pt}}\vspace{2.5mm}

\input{table_rows}

\vspace{4mm}
\begin{keybox}
{\sffamily\bfseries\footnotesize\color{Deep}Provenance and method}\\[1.2mm]
{\fontsize{7.2}{8.8}\selectfont
Period \textbf{27 September -- 3 October 2026}, Sunday to Saturday, \textbf{seven daily issues},
\textbf{125 records}, \textbf{125 unique DOIs}, no record without a DOI. This rollup is
\textbf{aggregation only}: it re-reads \texttt{state/corpus/*.json}, \texttt{state/metrics.csv} and
\texttt{state/trials.json} over the period and \textbf{queries no literature API}; every record was
gated by \texttt{check\_dois.py} at daily-issue time. Tier mix: \textbf{16 A, 55 B, 54 C}. The W40
weekly was due on the 3 October run. That run built the 2 October daily and deferred both; the runs of
6, 7, 8 and 9 October then harvested without authoring, three of them halting mid-write. The 3 October
daily and this rollup were built on the \textbf{9 October recovery run}, which also authored the
4--8 October dailies. The rollup includes the 3 October issue. One record admitted to the 3 October
issue (\texttt{10.1021/acs.est.6c07515}) could not be authored on that run and is \textbf{held, not
rejected}; it is absent from this pooled corpus and will enter a later issue.
Effect estimates are transcribed verbatim from the daily issues on heterogeneous contrasts and are
\textbf{not} meta-analysed. Abstracts remain the copyright of their respective publishers.\par}
\end{keybox}

\end{document}
